% dynamics/Vienna_reinterpretation-DE.tex
\chapter{Vienna-Funktionen als begrenztes Beispiel}
\label{chap:vienna-reinterpretation}

\section*{Evidenz vor Interpretation}

Der Name \emph{Vienna} kommt in historischem Projektmaterial und in
Implementierungen von Übergangsfunktionen vor. Diese Evidenz trägt die
Untersuchung bestimmter normierter Funktionsformen. Sie begründet weder eine
autoritative Krümmungs--Überhöhungs-Kopplung noch vollständige
Übergangsfamilien, validierte moderne Fahrzeugantwort oder allgemeine
gemeinsame Optimierung.

\begin{principle}[Prinzip der Vienna-Reinterpretation]
\label{princ:vienna-reinterpretation}
Eine Vienna-bezogene Funktion ist in dieser Thesis ein identifiziertes
historisches oder implementiertes Beispiel, dessen Eingaben, Ausgabegesetz
und Grenzen angegeben werden müssen.
\end{principle}

\begin{principle}[Kein-Spezialkurven-Prinzip]
\label{princ:vienna-no-special-curve}
Keine Vienna-bezogene Funktion wird hier zu einer kanonischen
AIM-Konstruktion erhoben.
\end{principle}

\section{Getrennte Gesetze}

Seien \(w=s/L\in[0,1]\) sowie \(f_\kappa(w)\) und \(f_u(w)\) zwei erklärte
normierte Funktionen. Eine illustrative Spezialisierung ist
\[
\kappa(s)=\kappa_0+\Delta\kappa\,f_\kappa(s/L),
\qquad
u(s)=u_0+\Delta u\,f_u(s/L).
\]
Diese Form ist mathematisch bestimmt, sobald beide Funktionen und Bereiche
vorgegeben sind. Sie sagt weder \(f_\kappa=f_u\) noch eine Kopplung ihrer
Parameter oder Optimalität aus.

\begin{definition}[Vienna-artige Zustandsformung]
\label{def:vienna-state-shaping}
„Vienna-artige Zustandsformung“ bleibt nur als historische Projektabsicht
erhalten: Eine identifizierte Übergangsfunktion dient zur Untersuchung einer
nachgelagerten Modellausgabe. Dies ist keine Kernel-Definition.
\end{definition}

\begin{definition}[Halbwelleninterpretation]
\label{def:half-wave-interpretation}
Eine Halbwelle ist eine kompositionale Implementierungs- oder Modellform,
wenn die identifizierte Funktion sie tatsächlich verwendet. Sie wird nicht
als Grammatik aller Übergangsfamilien behauptet.
\end{definition}

\begin{definition}[Dynamische Zustandsformung]
\label{def:dynamic-state-shaping}
Für ein identifiziertes Modell \(M\) ist die Änderung eines erklärten Gesetzes
mit Vergleich des resultierenden \(x_M^{\mathrm{pred}}\) eine illustrative
Sensitivitätsstudie. Die Bezeichnung „Formung“ liefert weder Ziel noch
Validierung.
\end{definition}

\section{Kompaktes Beispiel alternativer Ziele}

Für dieselben Randwerte und dieselbe Länge \(L\) werden zwei zulässige Paare
\((f_\kappa,f_u)\) betrachtet. Kandidat \(A\) minimiert
\[
J_\kappa=\int_0^L \bigl(\kappa'(s)\bigr)^2\,ds,
\]
Kandidat \(B\) dagegen
\[
J_r=\int_0^{L/v} r_{\mathrm{bal}}(t)^2\,dt
\]
für die ausdrücklich angenommene konstante Testgeschwindigkeit \(v\) und
\(M_{\mathrm{red}}\). Zieleinheiten und Bedeutungen sind verschieden. Ohne
vorgegebenen Zweck, Skalierung oder ausdrückliche Aggregation dominiert kein
Kandidat. Keines der Ziele ist ein Komfort-, Verschleiß-, Kosten- oder
Energieziel.

\section{Audit des Aussagestatus}

\begin{center}
\small
\begin{tabular}{p{0.29\linewidth}p{0.23\linewidth}p{0.38\linewidth}}
\toprule
Aussage & Status & Behandlung\\
\midrule
implementierte Funktionsform & implementationsbelegt & bei Nutzung exakte Implementierung nennen\\
historische dynamische Absicht & historische Projektabsicht & als Provenienz, nicht Validierung\\
getrennte Krümmungs-/Überhöhungsgesetze & illustrativ & Funktionen und Kopplungsannahme angeben\\
universelle Vienna-Kopplung & unbelegt & entfernt\\
validierte Fahrzeugantwort & unbelegt & benötigt Modell und Evidenz\\
allgemeine gemeinsame Optimierung & unbelegt & offene Arbeit, keine aktuelle Fähigkeit\\
\bottomrule
\end{tabular}
\end{center}

Dieselbe Zurückhaltung gilt für die Realisierung. Eine glatte Sollfunktion
kann mit qualifizierter Realisierungs- oder Beobachtungsevidenz verglichen
werden; eine Koordinatentransformation ist jedoch keine Konstruktion und eine
gemessene Abweichung keine dynamische Ursache. Das nächste Kapitel macht
diesen Vergleichsvertrag ausdrücklich.

% Beibehaltene Anker für Verweise auf die frühere längere Diskussion.
\label{prop:spectral-smoothing}
\label{princ:realization-survival}
\label{def:runtime-smoothness}
